% Bagian: computability
% Bab: computability-theory
% Subbagian: reducibility

\documentclass[../../../include/open-logic-section]{subfiles}

\begin{document}

\olfileid[id]{cmp}{thy}{red}
\olsection{Redusibilitas}

\begin{explain}
Sekarang kita mengetahui bahwa sekurang-kurangnya ada satu himpunan, yaitu
$K_0$, yang terenumerasi secara komputabel tetapi tidak dapat dihitung. Jelas
bahwa terdapat himpunan lain dengan sifat tersebut. Metode reduksi menyediakan
cara yang ampuh untuk menunjukkan bahwa himpunan-himpunan lain mempunyai
sifat-sifat ini tanpa harus terus-menerus kembali ke prinsip pertama.

Secara umum, suatu ``reduksi'' himpunan $A$ ke himpunan~$B$ adalah metode untuk
mengubah jawaban atas pertanyaan apakah !!{element}s berada dalam~$B$ menjadi
jawaban atas pertanyaan apakah !!{element}s berada dalam~$A$. Kita akan
berfokus pada gagasan yang disebut ``redusibilitas banyak-ke-satu,'' tetapi
tersedia banyak gagasan redusibilitas lain dengan sifat yang beragam. Gagasan
redusibilitas juga penting dalam kajian kompleksitas komputasional, yang turut
mempertimbangkan efisiensi. Sebagai contoh, suatu himpunan disebut
``NP-lengkap'' jika himpunan itu berada dalam NP dan setiap masalah NP dapat
direduksi kepadanya dengan menggunakan gagasan reduksi yang serupa dengan yang
dijelaskan di bawah, tetapi dengan syarat tambahan bahwa reduksi dapat dihitung
dalam waktu polinomial.

Kita sebenarnya sudah menggunakan gagasan reduksi secara tersirat. Definisikan
himpunan~$K$ dengan
\[
K = \Setabs{x}{\cfind{x}(x) \fdefined},
\]
yakni $K = \Setabs{x}{x \in W_x}$. Bukti kita bahwa masalah penghentian tidak
dapat diputuskan (\olref[hlt]{thm:halting-problem}) paling langsung menunjukkan
bahwa $K$ tidak dapat dihitung. Ingat bahwa $K_0$ adalah himpunan
\[
K_0 = \Setabs{\tuple{e, x}}{\cfind{e}(x) \fdefined },
\]
yakni $K_0 = \Setabs{\tuple{e,x}}{x \in W_e}$. Mudah memperluas setiap bukti
bahwa $K$ tidak dapat dihitung menjadi bukti bahwa $K_0$ tidak dapat dihitung:
jika $K_0$ dapat dihitung, kita dapat memutuskan apakah !!a{element}~$x$
berada dalam $K$ cukup dengan menanyakan apakah pasangan $\tuple{x, x}$ berada
dalam $K_0$. Fungsi~$f$ yang memetakan $x$ ke $\tuple{x, x}$ merupakan contoh
suatu \emph{reduksi} $K$ ke $K_0$.
\end{explain}

\begin{defn}
Misalkan $A$ dan $B$ himpunan bilangan asli. Suatu fungsi dapat dihitung
$f\colon \Nat \to \Nat$ merupakan \emph{reduksi banyak-ke-satu} dari $A$ ke
$B$ jika dan hanya jika, untuk setiap bilangan asli~$x$,
\[
x \in A \quad \text{jika dan hanya jika} \quad f(x) \in B.
\]
Jika terdapat reduksi~$f$ semacam itu, kita mengatakan bahwa $A$
\emph{tereduksi banyak-ke-satu} ke~$B$, yang ditulis $A \leq_m B$. Jika $A$
tereduksi banyak-ke-satu ke $B$ dan sebaliknya, maka $A$ dan $B$ disebut
\emph{ekuivalen banyak-ke-satu}, yang ditulis $A \equiv_m B$.
\end{defn}

Jika fungsi $f$ dalam definisi di atas kebetulan injektif, $A$ disebut
\emph{tereduksi satu-ke-satu} ke~$B$. Sebagian besar reduksi yang dijelaskan di
bawah memenuhi syarat yang lebih kuat ini, tetapi kita tidak akan menggunakan
fakta tersebut.

\begin{digress}
Memang benar, meskipun sama sekali tidak jelas, bahwa redusibilitas
satu-ke-satu benar-benar merupakan syarat yang lebih kuat daripada
redusibilitas banyak-ke-satu. Dengan kata lain, terdapat himpunan tak berhingga
$A$ dan~$B$ sedemikian sehingga $A$ tereduksi banyak-ke-satu ke~$B$, tetapi
tidak tereduksi satu-ke-satu ke~$B$.
\end{digress}

\end{document}

